\subsection{Acquiring a reading for a certified model decision}
\label{sec:decision-acquisition}
The width objective does not directly minimize requests to a decision.
We therefore consider a separate, single-query acquisition rule. Its
probabilities describe a declared acquisition surrogate; the final
all-completions certificate continues to use the original feasible set.
Write the current paired loss contrast as
$D(y)=d+\sum_g\sum_{j\in\mathcal M_g}c_jy_j$.
For a query $i\in\mathcal M_g$, let $l_i$ and $u_i$ be the extreme
coefficients remaining in block $g$, and let $A_g,B_g$ be the sharp
lower and upper contributions of $d$ and the other blocks. Conditional
on the answer $x\in[0,E_g]$, the updated sharp endpoints are
\begin{align}
 L_i(x)&=A_g+E_gl_i+(c_i-l_i)x,\notag\\
 U_i(x)&=B_g+E_gu_i+(c_i-u_i)x.\label{eq:decision-query-endpoints}
\end{align}
The answers that do not certify a strict ranking form the interval
$J_i=\{x\in[0,E_g]:L_i(x)\leq0\leq U_i(x)\}$.
If it is nonempty, denote its endpoints by $a_i,b_i$.

\begin{proposition}[One-query certification probability]
\label{prop:decision-query-probability}
Suppose $E_g>0$, block $g$ has at least two unknown coordinates, and
the normalized unknown targets in block $g$ have a Dirichlet
acquisition surrogate with positive parameters $(\alpha_j)_{j\in\mathcal M_g}$.
Set $\beta_i=\sum_{j\in\mathcal M_g\setminus\{i\}}\alpha_j$.
For an initially unresolved nonconstant contrast, the probability
that query $i$ certifies a strict ranking over all feasible completions is
\begin{equation}
 P_i=1-I_{b_i/E_g}(\alpha_i,\beta_i)
       +I_{a_i/E_g}(\alpha_i,\beta_i),
 \label{eq:decision-query-probability}
\end{equation}
with $P_i=1$ when $J_i$ is empty; an identically-zero contrast is
a known tie, not a strict ranking. Here $I_z$ is the regularized incomplete
beta function. Maximizing $P_i$ is optimal for this one-query surrogate
utility.
\end{proposition}
\begin{proof}
Allocating each block's remaining mass to an extreme coefficient gives
\eqref{eq:decision-query-endpoints}. The intersection of the two closed affine
half-lines defining $J_i$ is an interval. The Dirichlet marginal
$x/E_g\sim\operatorname{Beta}(\alpha_i,\beta_i)$ gives its probability.
An equality holding on an interval is included in $J_i$; isolated
endpoints have zero probability under the continuous
surrogate. The maximizing rule follows from the stated one-stage utility.
\end{proof}

For the uniform-simplex variant, $\alpha_j=1$ and
$P_i=1-(1-a_i/E_g)^{m_g-1}+(1-b_i/E_g)^{m_g-1}$.
For the machine-learning-guided variant, let $p_i$ be the frozen Ridge
forecast at an original six-hour block's coordinate and define
\[
 p_i^+=\max(p_i,0),\quad
 f_g=10^{-8}\max\!\left(\frac{1}{6}\sum_{j\in I_g}p_j^+,10^{-12}\right),
 \quad q_i=\max(p_i^+,f_g),\quad
 \alpha_i=6\kappa\frac{q_i}{\sum_{j\in I_g}q_j}.
\]
These parameters are fixed for all six coordinates before recovery.
Independent block surrogates are used. Conditioning on an observed
coordinate removes its parameter while preserving all remaining parameters;
their concentration is not renormalized. The positive concentration
$\kappa$ is selected using development data and frozen before confirmation.
Ties in $P_i$ are broken by smaller worst remaining width, then coordinate
index. The score uses exact zero thresholds; stopping retains the separately
specified numerical certificate rule.

This policy need not minimize total stopping cost. A narrow exact result
assumes an initially unresolved, nonconstant contrast, one uncertain block
with a positive exact total, two or three unknown coordinates after free
deductions, and no uncertainty elsewhere. Use exact strict-sign
certification and treat an exactly known tie as terminal. Under the stated
continuous Dirichlet surrogate, the nonconstant contrast is exactly zero
with probability zero. After an unresolved first query, at most one further
query determines the contrast; with two unknown coordinates the first
query already determines it. The expected number of queries following
first choice $i$ is therefore $2-P_i$, and maximizing $P_i$ minimizes this
surrogate cost in this restricted setting. This is not a result for
positive-tolerance regret stopping or an exact optimality assertion about
the floating-point stopping implementation. With more unknown coordinates
or several uncertain blocks, continuation states can matter and the
argument no longer applies.

The distinction from width optimization is strict. For one block with
$E=1$, $d=0$, and $c=(-4,-3,0,2)$, the one-query worst remaining widths are
$(5,6,6,4)$, whereas uniform-simplex certification probabilities are
$(8/27,27/125,0,1/27)$. The unique width optimum queries coefficient $2$;
the unique probability optimum queries coefficient $-4$. This constructed
example establishes objective mismatch, not an estimated application gain.
It is realizable by nonnegative forecasts: use $a_i=1-c_i/4$,
$b_i=1+c_i/4$ for these four missing coordinates, then add an observed
fifth coordinate with $y_5=a_5=0$, $b_5=\sqrt5$ to cancel the fixed
loss offset. Dividing by five to obtain MSE leaves the query choices
and probabilities unchanged.
Decision-region acquisition and expected-cost sequential testing are
established ideas; the contribution here is their explicit
aggregate-constrained score and its evaluation with an unchanged
all-completions stopping certificate.

